% GENERATED FILE. Edit canonical lessons, then run npm run generate:books.
% canonical-source-hash: fnv1a64:849d7b1492012b87
% canonical-lessons: PT-C363-cidadao, PT-C363-cidadania, PT-C363-direito, PT-C363-liberdade, PT-C363-igualdade

\chapter{Cidadãos}
\label{ch:pt-cidadaos}
\hlchaptermodality{363}

\begin{chapteropening}
\textbf{By the end of this chapter:} \emph{I can talk about citizens, their rights and freedoms, and equality.}

Complete the last lesson of chapter 363: I can talk about citizens, their rights and freedoms, and equality.
\end{chapteropening}

\section[o cidadão]{o cidadão — the citizen}

\label{lesson:PT-C363-cidadao}



\begin{quote}
Before the new one: say the Portuguese for the retirement, then the Portuguese for retired.
\end{quote}

\subsection*{You'll want to know: o cidadão}
\textbf{o cidadão} — \textquotedblleft{}the citizen\textquotedblright{}.

\begin{quote}
Todos os cidadãos podem votar aos dezoito anos.
\end{quote}

\begin{quote}
\emph{All citizens can vote at eighteen.}
\end{quote}

Plural \emph{os cidadãos}; feminine \emph{a cidadã}. \emph{O cartão de cidadão} is the Portuguese identity card.

\begin{cousinweb}
\emph{cidadão} is inherited from Old Portuguese \emph{çidadão} (EtymDB), built on \emph{cidade} \textquotedblleft{}city\textquotedblright{}, from Latin \emph{cīvitās} \textquotedblleft{}citizenship, city\textquotedblright{}, on \emph{cīvis} \textquotedblleft{}citizen\textquotedblright{} (Lewis \& Short).
\end{cousinweb}

\begin{grammarlens}[title={What you've built}]
One more word for telling, explaining and describing at B1.
\end{grammarlens}

\subsection*{Your turn}
\begin{itemize}
\raggedright
  \item \emph{Say it:} \emph{o cidadão}
  \item \emph{Say it:} \emph{o cidadão}, once more
  \item \emph{Say it:} a sentence of your own with \emph{o cidadão}
  \item \emph{Recall:} say the Portuguese for the retirement, then the Portuguese for retired, then say \emph{o cidadão} again
\end{itemize}

\begin{tcolorbox}[breakable,colback=teal!4,colframe=teal!35!black,title={Before you move on}]
What does o cidadão mean? (\textquotedblleft{}The citizen\textquotedblright{}.) Say it once more.
\end{tcolorbox}
\section[a cidadania]{a cidadania — the citizenship}

\label{lesson:PT-C363-cidadania}



\begin{quote}
Before the new one: say the Portuguese for retired, then the Portuguese for the citizen.
\end{quote}

\subsection*{You'll want to know: a cidadania}
\textbf{a cidadania} — \textquotedblleft{}the citizenship\textquotedblright{}.

\begin{quote}
Ela pediu a cidadania portuguesa.
\end{quote}

\begin{quote}
\emph{She applied for Portuguese citizenship.}
\end{quote}

\emph{Cidadania} is also a school subject about civic life.

\begin{cousinweb}
\emph{cidadania} is built on \emph{cidadão} \textquotedblleft{}citizen\textquotedblright{}, on \emph{cidade} \textquotedblleft{}city\textquotedblright{}, which comes from Latin \emph{cīvitās} \textquotedblleft{}citizenship, city\textquotedblright{}, on \emph{cīvis} \textquotedblleft{}citizen\textquotedblright{} (Lewis \& Short).
\end{cousinweb}

\begin{grammarlens}[title={What you've built}]
One more word for telling, explaining and describing at B1.
\end{grammarlens}

\subsection*{Your turn}
\begin{itemize}
\raggedright
  \item \emph{Say it:} \emph{a cidadania}
  \item \emph{Say it:} \emph{a cidadania}, once more
  \item \emph{Say it:} a sentence of your own with \emph{a cidadania}
  \item \emph{Recall:} say the Portuguese for retired, then the Portuguese for the citizen, then say \emph{a cidadania} again
\end{itemize}

\begin{tcolorbox}[breakable,colback=teal!4,colframe=teal!35!black,title={Before you move on}]
What does a cidadania mean? (\textquotedblleft{}The citizenship\textquotedblright{}.) Say it once more.
\end{tcolorbox}
\section[o direito]{o direito — the right; law (as a subject)}

\label{lesson:PT-C363-direito}



\begin{quote}
Before the new one: say the Portuguese for the citizen, then the Portuguese for the citizenship.
\end{quote}

\subsection*{You'll want to know: o direito}
\textbf{o direito} — \textquotedblleft{}the right; law (as a subject)\textquotedblright{}.

\begin{quote}
Todas as crianças têm direito à educação.
\end{quote}

\begin{quote}
\emph{All children have a right to education.}
\end{quote}

\emph{Ter direito a} — to have a right to. \emph{Estudar Direito} — to study law. Also the adjective right: \emph{a mão direita}.

\begin{cousinweb}
\emph{direito} is inherited from Latin \emph{dīrēctus} \textquotedblleft{}straight\textquotedblright{}, the participle of \emph{dīrigere} \textquotedblleft{}to set straight\textquotedblright{} (Lewis \& Short), built on \emph{regere} \textquotedblleft{}to guide, to rule\textquotedblright{}; \emph{-ct-} became \emph{-it-} as in \emph{feito}.
\end{cousinweb}

\begin{grammarlens}[title={What you've built}]
One more word for telling, explaining and describing at B1.
\end{grammarlens}

\subsection*{Your turn}
\begin{itemize}
\raggedright
  \item \emph{Say it:} \emph{o direito}
  \item \emph{Say it:} \emph{o direito}, once more
  \item \emph{Say it:} a sentence of your own with \emph{o direito}
  \item \emph{Recall:} say the Portuguese for the citizen, then the Portuguese for the citizenship, then say \emph{o direito} again
\end{itemize}

\begin{tcolorbox}[breakable,colback=teal!4,colframe=teal!35!black,title={Before you move on}]
What does o direito mean? (\textquotedblleft{}The right; law (as a subject)\textquotedblright{}.) Say it once more.
\end{tcolorbox}
\section[a liberdade]{a liberdade — the freedom, the liberty}

\label{lesson:PT-C363-liberdade}



\begin{quote}
Before the new one: say the Portuguese for the citizenship, then the Portuguese for the right.
\end{quote}

\subsection*{You'll want to know: a liberdade}
\textbf{a liberdade} — \textquotedblleft{}the freedom, the liberty\textquotedblright{}.

\begin{quote}
A liberdade de expressão é muito importante.
\end{quote}

\begin{quote}
\emph{Freedom of speech is very important.}
\end{quote}

\emph{Liberdade de expressão} — freedom of speech.

\begin{cousinweb}
\emph{liberdade} is inherited from Latin \emph{lībertās} \textquotedblleft{}the condition of a free person, freedom\textquotedblright{} (Lewis \& Short), through Old Portuguese \emph{liberdade} (EtymDB), built on \emph{līber} \textquotedblleft{}free\textquotedblright{}.
\end{cousinweb}

\begin{grammarlens}[title={What you've built}]
One more word for telling, explaining and describing at B1.
\end{grammarlens}

\subsection*{Your turn}
\begin{itemize}
\raggedright
  \item \emph{Say it:} \emph{a liberdade}
  \item \emph{Say it:} \emph{a liberdade}, once more
  \item \emph{Say it:} a sentence of your own with \emph{a liberdade}
  \item \emph{Recall:} say the Portuguese for the citizenship, then the Portuguese for the right, then say \emph{a liberdade} again
\end{itemize}

\begin{tcolorbox}[breakable,colback=teal!4,colframe=teal!35!black,title={Before you move on}]
What does a liberdade mean? (\textquotedblleft{}The freedom, the liberty\textquotedblright{}.) Say it once more.
\end{tcolorbox}
\section[a igualdade]{a igualdade — the equality}

\label{lesson:PT-C363-igualdade}



\begin{quote}
Before the new one: say the Portuguese for the right, then the Portuguese for the freedom.
\end{quote}

\subsection*{You'll want to know: a igualdade}
\textbf{a igualdade} — \textquotedblleft{}the equality\textquotedblright{}.

\begin{quote}
Ela luta pela igualdade entre homens e mulheres.
\end{quote}

\begin{quote}
\emph{She fights for equality between men and women.}
\end{quote}

\emph{Igualdade de oportunidades} — equal opportunities.

\begin{cousinweb}
\emph{igualdade} comes from Latin \emph{aequālitās} \textquotedblleft{}equality\textquotedblright{} through Old Portuguese \emph{ygualdade} (EtymDB), built on \emph{aequālis} \textquotedblleft{}equal\textquotedblright{}, on \emph{aequus} \textquotedblleft{}level, even\textquotedblright{} (Lewis \& Short).
\end{cousinweb}

\begin{grammarlens}[title={What you've built}]
One more word for telling, explaining and describing at B1.
\end{grammarlens}

\subsection*{Your turn}
\begin{itemize}
\raggedright
  \item \emph{Say it:} \emph{a igualdade}
  \item \emph{Say it:} \emph{a igualdade}, once more
  \item \emph{Say it:} a sentence of your own with \emph{a igualdade}
  \item \emph{Recall:} say the Portuguese for the right, then the Portuguese for the freedom, then say \emph{a igualdade} again
\end{itemize}

\begin{tcolorbox}[breakable,colback=teal!4,colframe=teal!35!black,title={Before you move on}]
What does a igualdade mean? (\textquotedblleft{}The equality\textquotedblright{}.) Say it once more.
\end{tcolorbox}
